\documentclass[12pt,a4paper]{article}

\usepackage{amsmath}
\usepackage{amssymb}
\usepackage{graphicx}
\usepackage{hyperref}
\usepackage{geometry}
\usepackage{fancyhdr}
\usepackage{enumitem}
\usepackage{titlesec}
\usepackage{listings}
\usepackage{xcolor}

% Page setup
\geometry{margin=1in}
\setlength{\headheight}{15pt}
\pagestyle{fancy}
\fancyhf{}
\rhead{SteveROS KBot ROS2 Control Architecture}
\lhead{\thepage}
\cfoot{\thepage}

% Colors
\definecolor{darkblue}{RGB}{0,51,102}
\definecolor{lightgray}{RGB}{240,240,240}

% Section formatting
\titleformat{\section}
{\normalfont\Large\bfseries\color{darkblue}}{\thesection}{1em}{}

\titleformat{\subsection}
{\normalfont\large\bfseries}{\thesubsection}{1em}{}

% Custom commands
\newcommand{\code}[1]{\texttt{#1}}
\newcommand{\ros}{\code{ROS2}}
\newcommand{\rc}{\code{ros2\_control}}

\begin{document}

\begin{titlepage}
    \centering
    \vspace*{2cm}
    
    \Huge\textbf{SteveROS KBot}
    \vspace{0.5cm}
    
    \Huge\textbf{ROS2 Control Architecture}
    \vspace{1cm}
    
    \Large\textbf{Technical Design Document}
    \vspace{2cm}
    
    \large Version 1.0\\
    \large \today
    \vspace{3cm}
    
    \begin{abstract}
        This document describes the architectural design for implementing ros2\_control 
        framework on the SteveROS KBot 20-DOF humanoid robot with Robstride actuators. 
        The architecture follows industry-standard patterns established by production-grade 
        humanoid robots including PAL Robotics TALOS and Pollen Robotics Reachy 2.
    \end{abstract}
    
\end{titlepage}

\tableofcontents
\newpage

\section{Introduction}

\subsection{Purpose}
This document defines the software architecture for the ros2\_control integration 
on the SteveROS KBot humanoid robot. The architecture provides a standard, maintainable, 
and extensible framework for low-level motor control while following best practices 
from production humanoid robot systems.

\subsection{Scope}
The architecture covers:
\begin{itemize}
    \item Hardware interface implementation for Robstride CAN-based actuators
    \item Package structure and organization
    \item ros2\_control configuration
    \item Controller setup and deployment
\end{itemize}

\subsection{Design Goals}

\begin{enumerate}
    \item \textbf{Standard Compliance} - Follow official ros2\_control framework patterns
    \item \textbf{Simplicity} - No over-engineering; maintain clean, minimal structure
    \item \textbf{Maintainability} - Clear separation of concerns
    \item \textbf{Extensibility} - Support future enhancements (simulation, MPC/WBC)
    \item \textbf{Professional Quality} - Match industry standards from production robots
\end{enumerate}

\section{Architecture Overview}

\subsection{Philosophy}
The architecture follows the principle of \textbf{convention over configuration}, adopting 
patterns proven in production humanoid robots. The design prioritizes:

\begin{itemize}
    \item Using the official \rc framework rather than custom solutions
    \item Single unified hardware interface for all joints
    \item Standard CAN-based communication using Linux SocketCAN
    \item Clean package separation following ros2 conventions
\end{itemize}

This approach aligns with industry standards used by:
\begin{itemize}
    \item PAL Robotics TALOS (32-DOF industrial humanoid)
    \item Pollen Robotics Reachy 2 (commercial service robot)
    \item Five Ages humanoid arms (research platform)
\end{itemize}

These production systems demonstrate that \rc provides the necessary flexibility 
and performance for full-size humanoid robots while maintaining code portability 
and community support.

\section{Package Structure}

\subsection{Overview}
The ros2 workspace follows a modular package organization with clear separation 
of concerns:

\subsection{Packages}

The workspace consists of three packages:

\begin{enumerate}
    \item \textbf{steveros\_hardware} - Hardware interface implementation
    \item \textbf{steveros\_description} - Robot model (URDF/XACRO)
    \item \textbf{steveros\_bringup} - Launch files and configurations
\end{enumerate}

\begin{table}[h]
    \centering
    \begin{tabular}{|l|p{6cm}|}
        \hline
        \textbf{Package} & \textbf{Contents} \\
        \hline
        steveros\_hardware & C++ hardware interface, CAN bus, Robstride protocol \\
        \hline
        steveros\_description & URDF/XACRO, meshes, RViz config, launch files \\
        \hline
        steveros\_bringup & Controller configs, bringup launch files \\
        \hline
    \end{tabular}
    \caption{Package structure}
\end{table}

This three-package structure follows the canonical ros2\_control demo pattern and 
is consistent with industry practice (e.g., TALOS, Reachy 2).

\section{Hardware Interface Design}

\subsection{Approach}
The hardware interface implements the \rc \code{SystemInterface} class, which is the 
standard approach for multi-joint robotic systems. This provides:

\begin{itemize}
    \item Unified access to all 20 joint actuators
    \item Standardized lifecycle management (initialize, configure, activate, deactivate)
    \item Thread-safe read/write interfaces for real-time control
    \item Plugin-based architecture for seamless simulation/mock testing
\end{itemize}

This approach is used by TALOS, Reachy 2, and the official ros2\_control demos.

\subsection{Communication Layer}
The hardware interface communicates with Robstride actuators via CAN bus using:

\begin{itemize}
    \item \textbf{Linux SocketCAN} - Native kernel-level CAN support
    \item \textbf{Non-blocking I/O} - Using epoll for responsive control
    \item \textbf{Per-axis abstraction} - Each motor maintains independent state/command
\end{itemize}

This pattern is validated by the ODrive hardware interface (ros\_odrive project), 
which implements identical patterns for CAN-based motor control.

\subsection{Joint Model}
Each joint exposes:

\begin{itemize}
    \item \textbf{Command interface}: position (primary)
    \item \textbf{State interfaces}: position, velocity, effort
\end{itemize}

This provides adequate functionality for position-based trajectory control while 
supporting future velocity and torque interfaces if needed.

\section{ros2\_control Configuration}

\subsection{Hardware Definition}
The robot description includes a single \code{<ros2\_control>} block defining:

\begin{itemize}
    \item Hardware plugin: \code{steveros\_hardware/RobstrideHardware}
    \item CAN interface parameter: \code{can0}
    \item All 20 joint definitions with motor IDs and limits
\end{itemize}

This single-system approach follows TALOS's pattern and is the industry standard 
for humanoid robots.

\subsection{Update Rate}
The controller manager operates at 100 Hz, which provides:

\begin{itemize}
    \item Sufficient bandwidth for position-based control
    \item Compatibility with joint trajectory controller requirements
    \item Room for future rate increases (500 Hz used by Reachy 2)
\end{itemize}

\section{Controller Configuration}

\subsection{Controllers}
The system uses standard \rc controllers:

\begin{enumerate}
    \item \textbf{joint\_state\_broadcaster} - Publishes joint states to ROS topics
    \item \textbf{joint\_trajectory\_controller} - Four instances for limb control:
    \begin{itemize}
        \item Right arm controller
        \item Left arm controller  
        \item Right leg controller
        \item Left leg controller
    \end{itemize}
\end{enumerate}

This controller distribution matches the physical limb structure and enables 
independent limb control while maintaining coordination through the controller 
manager.

\subsection{Interface Support}
Each controller is configured with:

\begin{itemize}
    \item \textbf{Command interfaces}: position
    \item \textbf{State interfaces}: position, velocity
    \item \textbf{Open-loop control}: disabled
    \item \textbf{Integration in goals}: enabled
\end{itemize}

These settings follow the standard configuration from the ros2\_control demos 
and are appropriate for trajectory-based motion.

\section{Design Influences}

\subsection{Primary References}
The architecture is derived from analysis of the following production systems:

\begin{table}[h]
    \centering
    \begin{tabular}{|l|c|p{5cm}|}
        \hline
        \textbf{System} & \textbf{DOF} & \textbf{Key Influence} \\
        \hline
        PAL Robotics TALOS & 32 & Single SystemInterface, async thread \\
        \hline
        Pollen Robotics Reachy 2 & 21-27 & 500Hz rate, dual-bus concept \\
        \hline
        ODrive ros2\_control & 2 & CAN communication patterns \\
        \hline
        K-Scale KBot & 20 & Robstride motor configuration \\
        \hline
        ros2\_control demos & - & Package structure, lifecycle patterns \\
        \hline
    \end{tabular}
    \caption{Architecture influence sources}
\end{table}

\subsection{Why These Patterns}
Each pattern was selected based on:

\begin{itemize}
    \item \textbf{Production validation} - Used in deployed humanoid robots
    \item \textbf{Community support} - Follows official framework conventions
    \item \textbf{Extensibility} - Enables future enhancements (simulation, MPC, WBC)
    \item \textbf{Maintainability} - Clear patterns understood by ros2 developers
\end{itemize}

\section{Future Extensions}

\subsection{Simulation Support}
The architecture supports future simulation integration through:

\begin{itemize}
    \item \textbf{MuJoCo} - Via mujoco\_ros2\_control plugin
    \item \textbf{Gazebo} - Via gazebo\_ros2\_control plugin
\end{itemize}

These can be enabled by changing the hardware plugin in the URDF, with no changes 
to higher-level code.

\subsection{Advanced Control}
The hardware interface design supports future implementation of:

\begin{itemize}
    \item \textbf{MPC/WBC} - Via legged\_control framework integration
    \item \textbf{RL policies} - Via Isaac Lab deployment
    \item \textbf{Torque control} - Adding effort command interface
\end{itemize}

These extensions require no changes to the core architecture.

\section{Summary}

The SteveROS KBot ros2\_control architecture follows industry-standard patterns 
established by production humanoid robots. Key characteristics:

\begin{itemize}
    \item \textbf{Standard} - Uses official ros2\_control framework
    \item \textbf{Clean} - Three-package structure with clear separation
    \item \textbf{Scalable} - Single SystemInterface for all joints
    \item \textbf{Proven} - Patterns validated by TALOS, Reachy 2, and fiveages
    \item \textbf{Extensible} - Ready for simulation and advanced control
\end{itemize}

This architecture provides a professional foundation for the KBot humanoid robot 
while maintaining simplicity and avoiding unnecessary complexity.

\section{References}

\begin{thebibliography}{99}

    \bibitem{talos}
    PAL Robotics, ``TALOS Robot ROS2 Packages,''
    \url{https://github.com/pal-robotics/talos_robot}, 2025.

    \bibitem{reachy}
    Pollen Robotics, ``Reachy 2 Core ROS2 Workspace,''
    \url{https://github.com/pollen-robotics/reachy2_core}, 2025.

    \bibitem{ros2control}
    ros-controls, ``ros2\_control Framework,''
    \url{https://github.com/ros-controls/ros2_control}, 2025.

    \bibitem{ros2controldemos}
    ros-controls, ``ros2\_control Demos,''
    \url{https://github.com/ros-controls/ros2_control_demos}, 2025.

    \bibitem{rosodrive}
    ODrive Robotics, ``ros\_odrive: ROS2 ros2\_control Driver for ODrive,''
    \url{https://github.com/odriverobotics/ros_odrive}, 2025.

    \bibitem{kbot}
    K-Scale Labs, ``KBot Humanoid Robot,''
    \url{https://github.com/kscalelabs/kbot}, 2025.

    \bibitem{fiveages}
    Five Ages, ``arms\_ros2\_control,''
    \url{https://github.com/fiveages-sim/arms_ros2_control}, 2025.

    \bibitem{mujoco}
    ros-controls, ``mujoco\_ros2\_control,''
    \url{https://github.com/ros-controls/mujoco_ros2_control}, 2025.

\end{thebibliography}

\end{document}
